\chapter{Base de datos}

Se utiliza PostgreSQL gestionado en el repositorio
\href{https://github.com/Serhuela/infra-postgresql-terraform}{infra-postgresql-terraform}.
El aprovisionamiento de las bases de datos, los esquemas y los usuarios corresponde
a esa infraestructura; este proyecto consume los recursos que se le proporcionan.

\section{Bases de datos y esquemas}

\begin{projecttable}
\toprule
\rowcolor{TableHeader}
\textbf{Elemento} & \textbf{Desarrollo} & \textbf{Producción} \\
\midrule
Base de datos & \path{psql_ddbb_dev_fund_nav} & \path{psql_ddbb_pro_fund_nav} \\ \hline
Esquema & \texttt{nav} & \texttt{nav} \\
\bottomrule
\end{projecttable}

El esquema \texttt{nav} se crea fuera de este proyecto y está disponible antes de
desplegar la aplicación. La aplicación no crea ni elimina esquemas.

\section{Usuarios}

\begin{projecttable}
\toprule
\rowcolor{TableHeader}
\textbf{Uso previsto} & \textbf{Desarrollo} & \textbf{Producción} \\
\midrule
DDL & \path{psql_user_dev_fund_nav_deploy} & \path{psql_user_pro_fund_nav_deploy} \\ \hline
DML & \path{psql_user_dev_fund_nav_bot} & \path{psql_user_pro_fund_nav_bot} \\
\bottomrule
\end{projecttable}

La infraestructura proporciona ambos usuarios para cada entorno. El usuario DDL
puede crear, modificar y eliminar objetos dentro del esquema asignado, pero no
crear ni eliminar esquemas. El usuario DML puede consultar y modificar los datos
de las tablas del esquema asignado (\texttt{SELECT}, \texttt{INSERT},
\texttt{UPDATE} y \texttt{DELETE}).
